\chapter*{Sommario}
\addcontentsline{toc}{chapter}{Sommario}

Il 1\textsuperscript{o} gennaio 2026 il \acrfull{cbam} dell'Unione europea è
passato da una fase transitoria di sola rendicontazione a un regime definitivo,
nel quale gli importatori devono restituire certificati commisurati alle
\glspl{embeddedemissions} dei beni coperti. La presente tesi si chiede se
quel passaggio abbia ridiretto in misura misurabile gli approvvigionamenti
europei.

La strategia di identificazione sfrutta il fatto che l'ambito di applicazione
è definito da un elenco esplicito di codici della \acrshort{cn}. I prodotti
coperti vengono confrontati con codici immediatamente contigui che ricadono
appena al di fuori dell'allegato --- beni simili, acquirenti simili, analoga
esposizione ciclica, ma nessun obbligo di certificazione --- in un disegno a
\acrfull{did} costruito sui dati mensili delle importazioni Comext di Eurostat.
Vengono esaminati tre esiti: i volumi importati, la riallocazione delle quote
di importazione fra paesi di origine con diversa intensità emissiva e i
\glspl{unitvalue} come approssimazione della traslazione dei costi sui prezzi.

[Risultati da completare una volta eseguita la stima.] L'analisi è concepita
per distinguere una riallocazione effettiva delle emissioni dalla semplice
\gls{resourceshuffling}, e per riportare quanta parte di un eventuale effetto
misurato sopravviva all'analisi di sensibilità sui trend paralleli proposta da
\textcite{rambachan2023}.

Il contributo è triplice: fornisce le prime stime quasi sperimentali del
regime definitivo del \acrshort{cbam}; separa l'effetto del prelievo da quello
della contestuale eliminazione dell'\gls{freeallocation}, che il dibattito di
policy tende a confondere; e documenta una strategia di identificazione che
resta utilizzabile man mano che le fasi successive amplieranno l'ambito
merceologico.

\medskip
\noindent\textbf{Parole chiave:} \thesisKeywords
